\pdfoutput=1
% =====================================================================================
%  torchtsetlin — five-page summary report (the document to cite)
%  Build:  make summary      (pdflatex x2)
%  Shares ttwhitepaper.sty, figures/ and data/ with the white paper (main.tex).
% =====================================================================================
\documentclass[10.5pt,a4paper]{article}

\usepackage[
  a4paper,
  top=22mm, bottom=22mm, left=22mm, right=22mm,
  headheight=14pt, headsep=8mm, footskip=10mm
]{geometry}

\usepackage{ttwhitepaper}
\usepackage[
  colorlinks=true, linkcolor=ttrust, citecolor=ttrust, urlcolor=ttrust,
  pdftitle={torchtsetlin: Fast, Interpretable Tsetlin Machines for the PyTorch Ecosystem},
  pdfauthor={Vajira Thambawita},
  pdfsubject={Technical report},
  pdfkeywords={Tsetlin machine, PyTorch, GPU, interpretable machine learning, propositional logic}
]{hyperref}

\newcommand{\ttversion}{0.3.2}
\newcommand{\ttdate}{19 October 2026}
% Generated by data/regen_literature.py from docs/assets/literature/literature.db.
\newcommand{\litNpapers}{55}
\newcommand{\litNdatasets}{128}
\newcommand{\litNresults}{3{,}295}
\newcommand{\litNusages}{252}
\newcommand{\litNmetrics}{204}
\newcommand{\litNtm}{1{,}974}
\newcommand{\litNbaseline}{1{,}321}
\newcommand{\litNproposed}{1{,}596}
\newcommand{\litNranked}{3{,}289}
\newcommand{\litNflagged}{3}
\newcommand{\litNdownload}{95}
\newcommand{\litYearMin}{2018}
\newcommand{\litYearMax}{2026}
\newcommand{\litDbKiB}{1{,}520}


\begin{document}
\thispagestyle{ttplain}

% ------------------------------------------------------------------ header band
\begin{tikzpicture}[remember picture, overlay]
  \fill[ttorange] (current page.north west) rectangle ($(current page.north east)+(0,-4.3cm)$);
  \fill[ttdeep] ($(current page.north west)+(0,-4.3cm)$) rectangle ($(current page.north east)+(0,-4.45cm)$);
  \begin{scope}[shift={($(current page.north east)+(-3.55cm,-0.95cm)$)}, scale=0.42,
                every node/.style={inner sep=0pt}]
    \foreach \r in {0,...,4}{
      \foreach \c in {0,...,6}{
        \pgfmathsetmacro{\v}{mod(\r*7+\c*5+int(\r*\c/2),9)}
        \ifdim \v pt > 5.5pt
          \fill[white, rounded corners=0.4pt] (\c,-\r) rectangle ++(0.8,0.8);
        \else
          \ifdim \v pt > 3.0pt
            \fill[white, opacity=0.40, rounded corners=0.4pt] (\c,-\r) rectangle ++(0.8,0.8);
          \else
            \draw[white, opacity=0.28, line width=0.5pt, rounded corners=0.4pt] (\c,-\r) rectangle ++(0.8,0.8);
          \fi
        \fi}}
  \end{scope}
  % Header text at fixed positions inside the band, so nothing can collide.
  \node[anchor=north west, inner sep=0pt, text=white, font=\sffamily\scriptsize\bfseries]
    at ($(current page.north west)+(2.2cm,-0.95cm)$)
    {TECHNICAL REPORT \ \textbar\ \ {\color{ttsand}\mdseries Software version \ttversion\ \ \textbar\ \ \ttdate}};
  \node[anchor=north west, inner sep=0pt, text=white, font=\sffamily\bfseries\fontsize{24}{28}\selectfont]
    at ($(current page.north west)+(2.2cm,-1.55cm)$) {torchtsetlin};
  \node[anchor=north west, inner sep=0pt, text=white, font=\sffamily\large]
    at ($(current page.north west)+(2.2cm,-2.6cm)$) {Fast, Interpretable Tsetlin Machines for the PyTorch Ecosystem};
  \node[anchor=north west, inner sep=0pt, text=white, font=\sffamily\small]
    at ($(current page.north west)+(2.2cm,-3.4cm)$)
    {\textbf{Vajira Thambawita}\ \ {\color{ttsand}Simula Metropolitan Center for Digital Engineering, Oslo, Norway\ \textbar\ \texttt{vajira@simula.no}}};
\end{tikzpicture}
\vspace*{2.55cm}

\begin{tikzpicture}
  \node[inner sep=0pt] (abs) {\begin{minipage}{0.955\linewidth}\small
    {\sffamily\bfseries\color{ttorange}Abstract.}\enspace
    The Tsetlin machine learns propositional rules with teams of Tsetlin automata driven by
    reward and penalty instead of gradients, which makes its models interpretable by
    construction and cheap to evaluate. Its software has grown up outside the deep-learning
    ecosystem, as C and CUDA kernels behind scikit-learn-style wrappers. \ttlibname expresses
    the whole family as ordinary \code{torch.nn.Module} objects: vanilla, weighted, coalesced,
    regression, convolutional over sequences, images and volumes, and dense per-pixel
    segmentation. The learnable state is a tensor of automaton counters. A model therefore
    moves across devices with \code{.to(device)}, serialises with \code{state\_dict()}, and
    learns from a mini-batch with one \code{model.update(x, y)} call. The central technical
    contribution is a batched formulation of Type~I\slash II feedback in which per-example
    automaton transitions become event counts realised as binomial draws, so there is one
    commit per batch. On an RTX~3090 this reaches $2.2\times10^{5}$ training examples\slash s
    and up to $287\times$ the same code on a 16-core CPU, while matching the sequential
    algorithm at moderate batch sizes. The package also ships an interactive, SQL-queryable
    catalog of \litNresults\ results on \litNdatasets\ datasets from \litNpapers\
    Tsetlin-machine papers.
  \end{minipage}};
  \begin{scope}[on background layer]
    \draw[ttorange!35, line width=0.8pt, rounded corners=4pt, fill=ttsand]
      ($(abs.north west)+(-0.8em,0.7em)$) rectangle ($(abs.south east)+(0.8em,-0.7em)$);
  \end{scope}
\end{tikzpicture}

\section{Statement of need}

A Tsetlin machine~\ttcite{1} represents each class with conjunctive clauses over Boolean
literals and their negations. A clause such as \code{IF x0 AND NOT x1} votes for or against a
class, and the vote count is the prediction. Its learned state is a matrix of small integers,
one Tsetlin automaton per (clause, literal) pair, updated by counting reward and penalty
events. The model is legible and its inference is bit-level cheap. That is why the family has
grown convolutional~\ttcite{2}, weighted~\ttcite{3}, coalesced~\ttcite{4} and regression
variants, and is used in clinical decision support and in hardware accelerators.

The reference software does not fit a modern research workflow. Each variant is a separate C
or CUDA kernel behind a scikit-learn-style class~\ttcite{9}. Models do not move between devices
like the rest of a PyTorch stack, they do not consume a \code{DataLoader} or serialise like
other models, and extending one means writing a kernel. \ttlibname starts from the opposite
end: if the automaton counters are a \code{torch} buffer, every variant is a few tensor
operations, and the PyTorch ecosystem~\ttcite{7} comes for free.

\section{Design}

\paragraph{State as tensors.} The automata live in an integer buffer \code{ta\_state} of shape
$(C, 2F)$; a literal is included when its counter crosses the action boundary $N$. The buffers
are not parameters (there are no gradients), so \code{to()}, \code{state\_dict()} and
\code{load\_state\_dict()} work unchanged. Learning is \code{update(x, y)}, the analogue of
\code{loss.backward(); optimizer.step()}. It returns the pre-update vote sums, so a trainer
logs batch metrics without a second forward pass.

\paragraph{One mechanism, many input structures.} Concrete models supply only an output layer
and a feedback policy through hooks on a shared base class. A different input structure
overrides the encoding, the clause evaluation and the feedback counts. Convolution over 1-D
sequences, 2-D images and 3-D volumes is one mixin; a clause fires if it matches any patch,
with thermometer-encoded position literals. Dense segmentation is a flat machine with one
example per pixel. Neither re-implements learning.

\paragraph{Batched feedback.} Classically, every example triggers a stochastic pass over all
$C\times 2F$ automata. For a batch of $B$ examples, \ttlibname instead forms the Type~I and
Type~II selection matrices $S_{\mathrm{I}}, S_{\mathrm{II}} \in \{0,1\}^{B\times C}$ and the
literal matrix $L \in \{0,1\}^{B\times 2F}$, and counts the events with two matrix products:
\begin{equation*}
  n^{\text{true}} = S_{\mathrm{I}}^{\!\top} L,\qquad
  n^{\text{false}} = S_{\mathrm{I}}^{\!\top} (1-L),\qquad
  n^{\text{II}} = S_{\mathrm{II}}^{\!\top} (1-L),
\end{equation*}
plus a column sum $n^{\text{Ib}}$ for non-firing clauses. The transition is then drawn in
closed form,
\begin{equation*}
  \Delta = \mathrm{Bin}\!\left(n^{\text{true}}, p_{\text{mem}}\right)
  - \mathrm{Bin}\!\left(n^{\text{false}} + n^{\text{Ib}}, \tfrac{1}{s}\right)
  + \min\!\left(n^{\text{II}}, N-\theta\right),
\end{equation*}
and committed once (Figure~\ref{fig:s-batched}). At $B=1$ this is the classical algorithm
term by term. For $B>1$ it approximates it the way mini-batch SGD approximates SGD, and
\code{feedback\_mode="sequential"} keeps the exact per-example path. A clause-size budget is
a hard cap~\ttcite{6}: newly included literals beyond it are reverted at commit.

\begin{figure}[htbp]
\centering
\resizebox{0.92\linewidth}{!}{% Figure: sequential per-example commits vs the batched event-counting formulation.
\begin{tikzpicture}[font=\sffamily\small]

% ===================================================================== A: sequential
\node[anchor=north west, font=\sffamily\bfseries\small, text=ttslate] at (0,0.75)
  {A\ \ Classical sequential algorithm \textemdash\ $B$ stochastic commits per mini-batch};

\foreach \i/\lbl in {0/{$x_1$},1/{$x_2$},2/{$x_3$}}{
  \pgfmathsetmacro{\x}{\i*4.45}
  \node[ttnode slate, anchor=north west, text width=3.7cm, align=center] at (\x,0)
    {\lbl\ \textemdash\ evaluate $\to$ select\\[2pt]
     \textbf{draw} $(C,2F)$ randoms\\[1pt]
     \textbf{write} the automata};
}
\node[anchor=north west, font=\sffamily, text=ttgray] at (13.5,-0.55) {$\cdots$};

\draw[ttarrow, draw=ttslate] (4.05,-0.65) -- (4.5,-0.65);
\draw[ttarrow, draw=ttslate] (8.5,-0.65) -- (8.95,-0.65);
\draw[ttarrow, draw=ttslate] (12.95,-0.65) -- (13.4,-0.65);

\node[ttlabel, anchor=north west, text width=14.3cm, font=\sffamily\small] at (0,-1.75)
  {Each example sees the automata the previous example left behind. That serial dependency is
   the \emph{point} of the algorithm --- and the reason a GPU sits idle: the work between two
   state writes is far too small to fill it.};

% ===================================================================== B: batched
\node[anchor=north west, font=\sffamily\bfseries\small, text=ttdeep] at (0,-3.35)
  {B\ \ Batched formulation \textemdash\ count the events, then commit \emph{once}};

\node[ttnode, anchor=north west, text width=1.85cm, align=center] (xb) at (0,-4.05)
  {$x_{1..B}$\\[1pt]\scriptsize $(B,F)$};
\node[ttnode, anchor=north west, text width=2.25cm, align=center] (lit) at (2.45,-4.05)
  {literals\\[1pt]\scriptsize $(B,2F)$};
\node[ttnode, anchor=north west, text width=2.65cm, align=center] (cl) at (5.35,-4.05)
  {clause outputs\\[1pt]\scriptsize $(B,C)$, one matmul};
\node[ttnode, anchor=north west, text width=2.45cm, align=center] (vt) at (8.65,-4.05)
  {vote sums\\[1pt]\scriptsize $(B,K)$, clipped};
\node[ttnode, fill=ttamber!25, draw=ttamber, anchor=north west, text width=2.85cm, align=center]
  (sl) at (11.75,-4.05) {feedback masks\\[1pt]\scriptsize $S_{\mathrm{I}}, S_{\mathrm{II}} \in \{0,1\}^{B\times C}$};

\draw[ttarrow] (2.0,-4.6) -- (2.4,-4.6);
\draw[ttarrow] (4.9,-4.6) -- (5.3,-4.6);
\draw[ttarrow] (8.2,-4.6) -- (8.6,-4.6);
\draw[ttarrow] (11.3,-4.6) -- (11.7,-4.6);

% ---- counting
\node[ttnode, anchor=north west, text width=13.9cm, align=center] (cnt) at (0.35,-5.95)
  {\textbf{event counts} \textemdash\ two matrix products over the chunk\\[4pt]
   $n^{\text{true}} = S_{\mathrm{I}}^{\!\top} L$ \qquad
   $n^{\text{false}} = S_{\mathrm{I}}^{\!\top}(1-L)$ \qquad
   $n^{\text{Ib}}_j = \textstyle\sum_b S^{\text{nofire}}_{\mathrm{I}}[b,j]$ \qquad
   $n^{\text{II}} = S_{\mathrm{II}}^{\!\top}(1-L)$};
\draw[ttarrow] (13.2,-5.25) .. controls (14.2,-5.4) and (14.2,-5.6) .. (13.2,-5.9);

% ---- commit
\node[ttnode dark, anchor=north west, text width=13.9cm, align=center] (com) at (0.35,-7.85)
  {\textbf{one commit} \textemdash\ the counts become automaton transitions\\[5pt]
   $\Delta \;=\; \mathrm{Bin}\!\left(n^{\text{true}},\, p_{\text{mem}}\right)
    \;-\; \mathrm{Bin}\!\left(n^{\text{false}} + n^{\text{Ib}},\, \tfrac{1}{s}\right)
    \;+\; \min\!\left(n^{\text{II}},\, N - \theta\right)$\\[5pt]
   $\theta \;\leftarrow\; \mathrm{clip}\!\left(\theta + \Delta,\; 0,\; 2N-1\right)$};
\draw[ttarrow, line width=1.1pt] (7.3,-7.25) -- (7.3,-7.8);

\node[ttlabel, anchor=north west, text width=14.3cm, font=\sffamily\small] at (0,-10.35)
  {Every stochastic decision of the sequential algorithm is an independent Bernoulli trial, so
   the trials one automaton receives in a batch sum to a \emph{binomial} draw --- one
   $(C,2F)$ random tensor per batch instead of one per example. At $B=1$ this reduces exactly
   to the classical algorithm.};
\end{tikzpicture}
}
\caption{The sequential algorithm (A) makes one stochastic $(C,2F)$ pass per example. The
  batched algorithm (B) accumulates the same events as counts and realises them in one
  binomial draw per batch, so the cost of an update is almost independent of $B$.}
\label{fig:s-batched}
\end{figure}

\begin{center}
\small\sffamily\sloppy
\captionof{table}{What ships in \ttlibname\ \ttversion.}
\label{tab:s-glance}
\begin{tabularx}{\linewidth}{@{}L{2.3cm}>{\raggedright\arraybackslash}X@{}}
\ttoprule
\thead{Models} & \clsname{TsetlinMachine} (optional integer clause weights), \clsname{CoalescedTsetlinMachine}
  (shared clause pool, multi-class or multi-label), \clsname{RegressionTsetlinMachine}; convolutional
  1-D, 2-D and 3-D variants of each; \clsname{SegmentationTsetlinMachine} and its coalesced form\\
\ttmidrule
\thead{Learning} & vote margin $T$, specificity $s$, memory depth $N$, boosted true-positive
  feedback, hard clause-size budget, drop-clause\slash drop-literal, focused negative sampling,
  batched or exact sequential feedback\\
\ttmidrule
\thead{Workflow} & Booleanization encoders (thermometer, one-hot, bit-plane, adaptive
  threshold, colour thermometer, hypervector); a Keras-style \clsname{Trainer} with callbacks;
  classification, regression and segmentation metrics; rule extraction and closed-form
  global\slash local importance~\ttcite{5}; plots\\
\ttmidrule
\thead{Footprint} & pure Python on PyTorch $\ge$ 2.0 and NumPy; no compiler, CUDA toolkit or C
  extension; Python 3.9--3.12; MIT licence\\
\ttbotrule
\end{tabularx}
\end{center}

\section{Usage}

Learning is a loop over \code{update}, or a Keras-style \clsname{Trainer} with callbacks. On
Noisy XOR with 40\% label noise, the model recovers the generating rule and prints it:

\begin{lstlisting}[style=ttpython]
import torchtsetlin as tt

x_train, y_train = tt.data.make_noisy_xor(5000, noise=0.4)    # Boolean features, int labels
x_test,  y_test  = tt.data.make_noisy_xor(2000, noise=0.0, seed=1)
model = tt.TsetlinMachine(n_features=12, n_classes=2, n_clauses=20, T=15, s=3.9).to("cuda")

trainer = tt.Trainer(model, batch_size=10, callbacks=[tt.EarlyStopping(patience=20)])
trainer.fit((x_train, y_train), epochs=100, val_data=(x_test, y_test))
print(trainer.evaluate((x_test, y_test)))      # {'accuracy': 0.99}
print(model.rules()[:2])                       # ['IF x0 AND NOT x1 THEN 1', 'IF x1 AND NOT x0 THEN 1']
\end{lstlisting}

Images, sequences and volumes use the convolutional classes with a \code{patch\_size}.
Continuous columns pass through a thermometer encoder whose feature names flow into the
rules, so a rule reads \code{worst perimeter>=105.9} rather than \code{x37}.
\code{model.eval()} matters: an empty clause votes while learning and abstains when
predicting.

\section{Performance}

An \code{update()} call touches every automaton, so its cost depends on the size of the
machine far more than on the batch. On an RTX~3090, a 784-feature machine with 500 clauses per
class trains at 427 examples\slash s at batch~1 and 220{,}028 at batch~4096
(Figure~\ref{fig:s-bench}). At batch~256 it runs at 87{,}869 examples\slash s against 406 on a
16-core CPU, a factor of $216\times$, and inference reaches $1.06\times10^{6}$ examples\slash
s. The advantage grows with the model, up to $287\times$ for the clause budgets the
image literature uses: a 113-second MNIST epoch becomes a 1.5-second one.

\begin{figure}[htbp]
\centering
% Figure: throughput vs batch size, training and inference, CPU vs GPU (MNIST-scale machine).
\pgfplotstableread[col sep=comma]{data/batch-mnist.csv}\batchmnist
\pgfplotstableread[col sep=comma]{data/batch-xor.csv}\batchxor

\begin{tikzpicture}
\begin{groupplot}[
    ttaxis,
    group style={group size=2 by 1, horizontal sep=1.45cm},
    width=0.5\linewidth, height=5.6cm,
    xmode=log, ymode=log,
    log basis x=2,
    xlabel={batch size},
    xtick={1,4,16,64,256,1024,4096},
    xticklabels={1,4,16,64,256,1k,4k},
    legend pos=north west,
  ]

  % ---------------------------------------------------------------- training
  \nextgroupplot[title={Training}, ylabel={examples / s}]
    \addplot[ttgpu] table[x=batch, y=gpu_train] {\batchmnist};
    \addlegendentry{GPU (RTX 3090)}
    \addplot[ttcpu] table[x=batch, y=cpu_train] {\batchmnist};
    \addlegendentry{CPU (16 threads)}
    \addplot[ttthird, dashed, mark=none] table[x=batch, y=gpu_train] {\batchxor};
    \addlegendentry{tiny machine, GPU}
    \addplot[color=ttgray, dotted, line width=1.1pt, mark=none]
      table[x=batch, y=cpu_train] {\batchxor};
    \addlegendentry{tiny machine, CPU}

  % ---------------------------------------------------------------- inference
  \nextgroupplot[title={Inference}, ylabel={examples / s}]
    \addplot[ttgpu] table[x=batch, y=gpu_infer] {\batchmnist};
    \addlegendentry{GPU (RTX 3090)}
    \addplot[ttcpu] table[x=batch, y=cpu_infer] {\batchmnist};
    \addlegendentry{CPU (16 threads)}
\end{groupplot}
\end{tikzpicture}

\caption{Throughput against batch size, GPU and CPU, for a 784-feature, 500-clause-per-class
  machine (log--log). The dashed line is a 12-feature, 10-clause machine: less work than a
  kernel launch, so the GPU only draws level at batch 1024.}
\label{fig:s-bench}
\end{figure}

Batching costs little accuracy when the batch is moderate. On Noisy XOR, the benchmark of the
original paper, batches up to about 50 land in the same accuracy band as the exact sequential
algorithm, and batch 200 degrades clearly (Figure~\ref{fig:s-fidelity}).

\begin{figure}[htbp]
\centering
% Figure: what batching costs in accuracy — Noisy XOR, 3 seeds per batch size.
% A strip plot rather than bars: the y axis does not start at zero, and bar lengths on a
% truncated axis would overstate the differences.
\begin{tikzpicture}
\begin{axis}[
    ttaxis,
    width=0.66\linewidth, height=5.6cm,
    symbolic x coords={1, 10, 50, 200},
    xtick=data,
    xlabel={batch size \ (1 = the exact sequential algorithm)},
    ylabel={test accuracy},
    ymin=0.73, ymax=1.03,
    ytick={0.75,0.80,0.85,0.90,0.95,1.00},
    yticklabel style={/pgf/number format/fixed, /pgf/number format/zerofill,
                      /pgf/number format/precision=2},
    enlarge x limits=0.18,
    legend style={at={(0.5,-0.32)}, anchor=north, legend columns=3},
  ]

  % seed-to-seed range
  \draw[ttgray!45, line width=1.2pt] (axis cs:1,0.954)   -- (axis cs:1,0.993);
  \draw[ttgray!45, line width=1.2pt] (axis cs:10,0.931)  -- (axis cs:10,0.993);
  \draw[ttgray!45, line width=1.2pt] (axis cs:50,0.877)  -- (axis cs:50,0.993);
  \draw[ttgray!45, line width=1.2pt] (axis cs:200,0.771) -- (axis cs:200,0.889);

  \addplot[only marks, mark=o, mark size=2.6pt, draw=ttslate, line width=0.9pt]
    coordinates {(1,0.954) (1,0.955) (1,0.993)
                 (10,0.942) (10,0.993) (10,0.931)
                 (50,0.993) (50,0.965) (50,0.877)
                 (200,0.788) (200,0.771) (200,0.889)};
  \addlegendentry{individual seeds}

  \addplot[only marks, mark=diamond*, mark size=3.6pt, fill=ttorange, draw=ttdeep,
           line width=0.7pt, point meta=explicit symbolic, nodes near coords,
           nodes near coords style={font=\sffamily\scriptsize, color=ttdeep,
                                    anchor=west, xshift=5pt}]
    coordinates {(10,0.955) [0.955] (50,0.945) [0.945] (200,0.816) [0.816]};
  \addlegendentry{mean of 3 seeds}

  % batch 1 is labelled below-right so the label clears the reference line it defines
  \addplot[only marks, mark=diamond*, mark size=3.6pt, fill=ttorange, draw=ttdeep,
           line width=0.7pt, point meta=explicit symbolic, nodes near coords, forget plot,
           nodes near coords style={font=\sffamily\scriptsize, color=ttdeep,
                                    anchor=north west, xshift=3pt, yshift=-2pt}]
    coordinates {(1,0.967) [0.967]};

  \addplot[ttamber, dashed, line width=1.0pt, mark=none, forget plot]
    coordinates {(1,0.967) (200,0.967)};
  \addlegendimage{ttamber, dashed, line width=1.0pt, mark=none}
  \addlegendentry{sequential mean}
\end{axis}
\end{tikzpicture}

\caption{Noisy XOR test accuracy against batch size, 3 seeds; batch 1 is the exact sequential
  algorithm.}
\label{fig:s-fidelity}
\end{figure}

\noindent The library therefore recommends batches of 16--64
for flat models and 8--32 for convolutional ones. Every number above comes from the benchmark
driver in the repository, which regenerates the documentation's CPU-versus-GPU page from its
own records, so it can be re-run on other hardware with one command.

\section{An interactive review of the literature}

The package ships a catalog of every dataset and reported number in \litNpapers\
Tsetlin-machine papers (\litYearMin--\litYearMax): \litNdatasets\ datasets, \litNusages\
dataset usages with the splits as reported, and \litNresults\ results, covering both the
proposed models and their baselines. One JSON record per paper is the source of truth. It is
compiled into a \litDbKiB~KiB SQLite database, which the documentation opens in the browser
with \code{sql.js}. Readers filter, group and rank results there, see the best value per
group, compare the best Tsetlin-machine result with the best baseline per dataset, and run
their own SQL. Values are ranked only within the same metric, unit and split, and a metric's
direction decides what counts as best. Papers with a known reliability problem carry a
caution flag (\litNflagged\ papers). Figure~\ref{fig:s-literature} shows two of its views.

\begin{figure}[htbp]
\centering
\resizebox{\linewidth}{!}{% Figure: the literature catalog. Data: data/lit-top.csv and data/lit-best.csv, both written by
% data/regen_literature.py from docs/assets/literature/literature.db.
\pgfplotstableread[col sep=comma]{data/lit-top.csv}\littop
\pgfplotstableread[col sep=comma]{data/lit-best.csv}\litbest
\begin{tikzpicture}
\begin{axis}[
    ttaxis,
    name=top,
    width=0.47\linewidth, height=6.4cm,
    xbar, bar width=6pt,
    xmin=0,
    xlabel={papers using the dataset},
    ytick=data, y dir=reverse,
    yticklabels from table={\littop}{dataset},
    yticklabel style={font=\ttfamily\scriptsize, color=ttink},
    ymajorgrids=false,
    enlarge y limits=0.06,
    nodes near coords,
    nodes near coords style={font=\sffamily\tiny, color=ttdeep},
    title={Most-used datasets},
  ]
  \addplot[ttbar] table[x=papers, y expr=\coordindex] {\littop};
\end{axis}
\begin{axis}[
    ttaxis,
    at={($(top.east)+(1.9cm,0)$)}, anchor=west,
    width=0.53\linewidth, height=6.4cm,
    ybar, bar width=6.5pt,
    ymin=75, ymax=101,
    ylabel={best reported test accuracy (\%)},
    xtick=data,
    xticklabels from table={\litbest}{dataset},
    xticklabel style={font=\sffamily\scriptsize, color=ttink, rotate=35, anchor=north east},
    xmajorgrids=false,
    enlarge x limits=0.1,
    legend style={at={(0.5,1.02)}, anchor=south, legend columns=2},
    title={TM vs.\ the baselines TM papers report},
    title style={yshift=13pt},
  ]
  \addplot[ttbar] table[x expr=\coordindex, y=tm] {\litbest};
  \addlegendentry{best TM model}
  \addplot[ttbarcpu] table[x expr=\coordindex, y=baseline] {\litbest};
  \addlegendentry{best non-TM baseline}
\end{axis}
\end{tikzpicture}
}
\caption{From the literature catalog. \emph{Left:} the datasets used by the most papers.
  \emph{Right:} the best reported Tsetlin-machine test accuracy next to the best non-TM
  baseline reported in the same papers, excluding flagged papers and federated per-client
  accuracies. All values are as reported by their authors, not reproduced.}
\label{fig:s-literature}
\end{figure}

\section{Limitations}

\paragraph{Batched feedback is an approximation.} The default path trades some fidelity for
throughput, and the characterisation is empirical: there is no proof that the batched update
converges to the fixed point of the sequential one. The exact algorithm is one argument away.

\paragraph{Dense tensors, not bit-packed kernels.} Literals are floats in a matrix product, not
bits in a machine word, and clause sparsity is not exploited by indexing. On a small machine
on a CPU, the reference C implementations are faster; the advantage of \ttlibname appears with
model size and a GPU.

\paragraph{Coverage of the literature.} Graph Tsetlin machines, TM composites~\ttcite{8}, contracting
machines with absorbing automata and several other recent variants are not implemented yet.
Training is single-device and single-process.

\section{Availability and citation}

\begin{center}
\small\sffamily\sloppy
\begin{tabularx}{\linewidth}{@{}L{2.6cm}>{\raggedright\arraybackslash}X@{}}
\ttoprule
\thead{Install} & \code{pip install torchtsetlin} \ (version \ttversion, \ttdate)\\
\thead{Source} & \url{https://github.com/vlbthambawita/torchtsetlin} (MIT licence)\\
\thead{Documentation} & \url{https://vajira.info/torchtsetlin/}, including the literature
  catalog and the generated benchmark pages\\
\thead{Quality gates} & \code{ruff} and the full test suite on Python 3.9, 3.11 and 3.12 in CI;
  releases via PyPI Trusted Publishing after a clean-environment smoke import\\
\ttbotrule
\end{tabularx}
\end{center}

If you use \ttlibname, please cite this report, together with the Tsetlin-machine papers whose
algorithms you rely on:

\begin{lstlisting}[style=ttshell,basicstyle=\ttfamily\scriptsize\color{ttink}]
@techreport{thambawita2026torchtsetlin,
  author      = {Thambawita, Vajira},
  title       = {torchtsetlin: Fast, Interpretable Tsetlin Machines for the PyTorch Ecosystem},
  institution = {Simula Metropolitan Center for Digital Engineering},
  type        = {Technical report},
  year        = {2026},
  note        = {Software version 0.3.2},
  url         = {https://github.com/vlbthambawita/torchtsetlin}
}
\end{lstlisting}

\section*{References}
{\small\sloppy
\setlength{\parskip}{0.15em}
\begin{list}{\textcolor{ttorange}{[\arabic{ttrefcnt}]}}%
  {\usecounter{ttrefcnt}\setlength{\leftmargin}{2.1em}\setlength{\labelwidth}{1.7em}%
   \setlength{\labelsep}{0.4em}\setlength{\itemsep}{0.15em}\setlength{\topsep}{0.3em}}
\item O.-C. Granmo. \emph{The Tsetlin Machine --- A Game Theoretic Bandit Driven Approach to
  Optimal Pattern Recognition with Propositional Logic}. arXiv:1804.01508, 2018.
\item O.-C. Granmo, S. Glimsdal, L. Jiao, M. Goodwin, C. W. Omlin, G. T. Berge.
  \emph{The Convolutional Tsetlin Machine}. arXiv:1905.09688, 2019.
\item A. Phoulady, O.-C. Granmo, S. R. Gorji, H. A. Phoulady. \emph{The Weighted Tsetlin
  Machine: Compressed Representations with Weighted Clauses}. arXiv:1911.12607, 2019.
\item S. Glimsdal, O.-C. Granmo. \emph{Coalesced Multi-Output Tsetlin Machines with Clause
  Sharing}. arXiv:2108.07594, 2021.
\item C. D. Blakely, O.-C. Granmo. \emph{Closed-Form Expressions for Global and Local
  Interpretation of Tsetlin Machines}. arXiv:2007.13885, 2020.
\item K. D. Abeyrathna et al. \emph{Building Concise Logical Patterns by Constraining Tsetlin
  Machine Clause Size}. IJCAI 2023, arXiv:2301.08190.
\item A. Paszke et al. \emph{PyTorch: An Imperative Style, High-Performance Deep Learning
  Library}. NeurIPS 2019.
\item O.-C. Granmo. \emph{TMComposites: Plug-and-Play Collaboration Between Specialized
  Tsetlin Machines}. arXiv:2309.04801, 2023.
\item Centre for Artificial Intelligence Research (CAIR). \emph{TMU: the Tsetlin Machine
  Unified library}. \url{https://github.com/cair/tmu}.
\end{list}}
\end{document}
